\section{Introduction}

The project develops a one-way coupled tsunami simulation workflow for the 2011 Tohoku event, linking coseismic source generation, regional propagation and local three-dimensional hydrodynamic impact. The completed G6 gate establishes the theoretical 2D--3D hybrid model and its real-event integration baseline. It does not claim observational validation or physical calibration.

The motivation is methodological: regional nonlinear shallow-water models are efficient for source-to-coast propagation, while local URANS--VOF simulations resolve free-surface interaction with coastal defences and buildings. Prior hybrid and comparative studies support this separation of roles \autocite{fujima2002hybrid,takase2016hybrid,qin2018comparison}. The Tohoku event provides the source, observation and impact context \autocite{fujii2011tsunamisource,mori2012nationwide,fine2013japan}.

This document consolidates the Research workstream into the authoritative project report. The Research workstream remains the evidence layer; this report records the accepted G6 formulation, traceability, verification status and the methodology required before calibration and validation can be executed.
